\documentclass[10pt,a4paper]{amsart}
\usepackage[margin=19mm]{geometry}
\usepackage{amsmath,amssymb,mathtools}
\usepackage[scr=rsfs,cal=euler]{mathalfa}
\usepackage{xfrac}
\usepackage[unicode,hidelinks,bookmarks=false]{hyperref}
\newcommand{\ie}{i.e.\ }
\newcommand{\cf}{cf.\ }
\newcommand{\resp}{respectively\ }
\newcommand{\Subsection}{\subsection}
\providecommand{\nobreakdash}{\nobreak\hbox{-}}
\setlength{\emergencystretch}{2em}
\multlinegap0pt
\usepackage{bm}
\usepackage{bbm}
\usepackage{dsfont}
\usepackage{xspace}
\usepackage{cancel}
\usepackage{mathtools}
\usepackage{amsthm}
\makeatletter
\@ifundefined{dummy}{\newtheorem{dummy}{Dummy}}{}
\@ifundefined{lemma}{\newtheorem{lemma}[dummy]{Lemma}}{}
\makeatother
\newcommand{\N}{\ensuremath{\mathbb{N}}}
\newcommand{\R}{\ensuremath{\mathbb{R}}}
\newcommand{\diff}{\ensuremath{\operatorname{d}\!}}
\title[Radon-Nikodym derivative of inhomogeneous Brownian last pass]{Radon-Nikodym derivative of inhomogeneous Brownian last passage percolation}
\author{Pantelis Tassopoulos, Sourav Sarkar}
\date{Source: arXiv:2509.19414v3 (CC BY 4.0)}
\begin{document}
\maketitle

\noindent\emph{Source and licence.} Radon-Nikodym derivative of inhomogeneous Brownian last passage percolation. Version arXiv:2509.19414v3 (CC BY 4.0), distributed under the Creative Commons Attribution 4.0 International licence, as recorded in the arXiv licence metadata for this exact version and shown on the abstract page https://arxiv.org/abs/2509.19414v3. Full rights notice: licenses/arxiv-2509.19414-CC-BY-4.0.txt.\par\smallskip

\noindent\textit{Editorial scope.} Counted unit: the Lemma whose source label is \texttt{lemma: difference bound} in the article (source0.tex lines 1597--1606, source label \texttt{lemma: difference bound}), delivered together with the article's own complete original proof of it (source0.tex lines 1608--1693, one \texttt{proof} environment of 86 lines). No mathematics is written, repaired or re-derived: statement and proof are the article's own text line for line. Required context retained uncounted: lemma: integral estimate (lines 1793--1814); eq: poly-gaussian bounds (lines 1795--1797). Every definition, display and label of those lines is retained unchanged. Omitted from this package and not redistributed: 1--1596, 1607--1607, 1694--1792, 1798--2083. Nothing inside a retained excerpt is omitted or altered: every retained body is delivered exactly as the article's lines are written, so no omission receipt of any kind accompanies this package. Citations: the retained excerpts carry no citation command at all, so no bibliography entry of the article is transcribed and no reference is redistributed. Reference adaptation: none was needed and none was made; every reference inside the delivered unit resolves inside it, so no unresolved reference or citation is produced. Printed slips kept literal: no printed word, symbol, hypothesis or number of the counted statement or of its proof was altered. Layout and class: the article's own class is \texttt{amsart} with packages including fontenc, lmodern, mathtools, amsfonts, amssymb, mathrsfs, bm, stmaryrd, graphicx, tikz, pgfplots, pgfplotstable, float, svg; the wrapper is the lane's pinned \texttt{amsart} editorial wrapper at 10pt on A4, which restates the theorem-like environments the retained text needs and adds the article's own macro definitions that the retained text needs (\texttt{\textbackslash N}, \texttt{\textbackslash R}, \texttt{\textbackslash diff}), with extra wrapper packages (bm, bbm, dsfont, xspace, cancel, mathtools). The delivered numbering is the unit's own. Assets: no retained span contains a figure environment, a graphics-inclusion command, a PDF-page inclusion, a table or a TikZ picture; no image file is copied into this package, the package declares no asset dependency and no image file is redistributed with it. Licence: version arXiv:2509.19414v3 of this article is distributed under the Creative Commons Attribution 4.0 International licence, as recorded in the arXiv licence metadata for this exact version and shown on the abstract page https://arxiv.org/abs/2509.19414v3. A full rights notice is delivered at licenses/arxiv-2509.19414-CC-BY-4.0.txt. Characters: every retained excerpt is pure ASCII, so the delivered engine, XeLaTeX with its default Latin Modern text font, reports no missing character.
\begin{lemma}\label{lemma: integral estimate}
Fix $m\in \N$, let $g:\R\to \R $ be a smooth function such that there exists an $N\in \N$ with the following estimate
\begin{equation*}\label{eq: poly-gaussian bounds}
    C_1(\theta, \beta)\langle y\rangle^{-1}\cdot \mathrm{e}^{-\beta(y)_-}\cdot \mathrm{e}^{-\theta(y)^2_-} \le g(y) \le C_2(N, \theta, \beta) \langle y\rangle^{N}\cdot \mathrm{e}^{-\theta(y)^2_-},\quad y\in \R.
\end{equation*}
for some $\theta >0$, $C_1(\theta, \beta), C_2(N, \theta, \beta)>0$. Then, we have the upper and lower bounds
    \begin{align*}
          &C^\prime_1(\theta, \beta)\langle (x)_-\lor 1 \rangle ^{-1} \mathrm{e}^{-(\theta + 1/2)(x)_-^2}\mathrm{e}^{-(x)_-\big(\beta+3\big)}\le \displaystyle\int^x_{-\infty}\mathrm{e}^{-\frac{(y)^2}{2}} g(y)(x-y)^m\diff y \\
         & \le C^\prime_2(N, m, \theta, \beta)\langle x\rangle ^{N+m+1}\mathrm{e}^{-(\theta+1/2)(x)_-^2},\quad x\in \R.
    \end{align*}
for $C^\prime_1(\theta, \beta) = C_1(\theta, \beta)\frac{c}{(2\theta+1)^m}\frac{2^{1/2}}{\frac{\beta+3}{2\theta +1} +  2^{1/2}}\mathrm{e}^{(\theta + 1/2)(\frac{\beta+2}{2\theta +1})^2}\mathrm{e}^{-(\theta + 1/2)\big(\frac{\beta+3}{2\theta +1}\big)^2}$ and\\ $C^\prime_2(N, m, \theta, \beta) =  d^{(m+N)\log (m+N)}\cdot C_2(N, \theta, \beta)$, for a universal constant $c>0$, that is of the same form as \ref{eq: poly-gaussian bounds}. Furthermore, for all $\lambda>0$, $K\in \N$
\begin{align*}
    \displaystyle\int^x_{-\infty} \langle y \rangle^K \mathrm{e}^{\lambda(y)_-}\mathrm{e}^{-\frac{y^2}{2}} g(y)(x-y)^m\diff y &\le C_{N,K, M, m, \theta, \beta, \lambda}  \langle x\rangle ^{N+K+m+1}\mathrm{e}^{\big(\beta +3\big)\cdot (x)_-}\\
    & \cdot\displaystyle\int^x_{-\infty}\mathrm{e}^{-\frac{(y)^2}{2}} g(y)(x-y)^m\diff y\, \quad x\in \R.
\end{align*}
where 
\begin{align*}
         & C_{N,K,M,m,\theta, \beta, \lambda}=C_1(\theta, \beta) \cdot C_2(N, \theta, \beta) d^{(N+K+m)\log(N+K+m)}\cdot(2\theta+1)^m  \\
         & \cdot \mathrm{e}^{\frac{\lambda^2}{2\theta + 1}} \mathrm{e}^{(\theta + 1/2)\big(\frac{\beta+3}{2\theta +1}\big)^2} \mathrm{e}^{-(\theta + 1/2)\big(\frac{\beta+2}{2\theta +1}\big)^2}
    \end{align*}
for a universal constant $d>0$ and $\langle\cdot \rangle = (\cdot ^2 + 1)^{1/2}$. 
\end{lemma}

\begin{lemma}\label{lemma: difference bound} Fix $k\in \N, \ell> 0$ and $m_1, m_2, \cdots, m_k, m_{k+1} \in \N$ and consider the function
\begin{align*}
g_{m_1, \cdots, m_k}(b) = \displaystyle\int_{x_1\le x_2\le \cdots\le x_k\le b}\prod_{i=1}^{k}\mathrm{e}^{-\frac{x_i^2}{2}}\cdot(x_2-x_1)^{m_1}\cdot \cdots \cdot (b-x_k)^{m_k}\diff x_1\cdots \diff x_{k}\,, b\in \R.
\end{align*}
Then there exists a universal constant $d>0$ such that 
    $$
     g_{m_1, \cdots, m_k}(b+1)\le C_{k, m_1, \cdots, m_k} \cdot (1+\mathrm{e}^{-kb})\cdot g_{m_1, \cdots, m_k}(b), \quad b\in \R
    $$
    where $C_{k+1, m_1, \cdots, m_{k+1}} = 2^{m_{k+1}}(C_{k,\underline{m}}\mathrm{e}^{k+1}\lor \mathrm{e}^{d\sum_{i=1}^k m_i\log m_i})$, and $C_{1,m_1} = O(2^{m_1})$.
\end{lemma}

\begin{proof}
    Indeed, we proceed by induction. To see this note that for the inductive step
    \begin{equation*}
    \begin{array}{cc}
         g_{m_1, \cdots, m_{k+1}}(b) &= \displaystyle\int^b_{-\infty}\mathrm{e}^{-\frac{(y)^2}{2}} g_{m_1, \cdots, m_{k}}(y)(b-y)^{m_{k+1}}\diff y\\
         & = \displaystyle\int^\infty_{0}\mathrm{e}^{-\frac{(b-y)^2}{2}} g_{m_1, \cdots, m_{k}}(b-y)(y)^{m_{k+1}}\diff y,\quad b\in \R
    \end{array}
    \end{equation*}
    and
\begin{align*}
& \frac{g_{m_1, \cdots, m_{k+1}}(b+1)}{g_{m_1, \cdots, m_{k+1}}(b)} =  \frac{
    \displaystyle \int_{-\infty}^b \mathrm{e}^{-\frac{y^2}{2}} g_{m_1, \cdots, m_k}(y)
    (b+1 - y)^{m_{k+1}} \, \diff y
}{
    \displaystyle \int_{-\infty}^b \mathrm{e}^{-\frac{y^2}{2}} g_{m_1, \cdots, m_k}(y)
    (b - y)^{m_{k+1}} \, \diff y
}
\\
& + \frac{
    \displaystyle \int_b^{b+1} \mathrm{e}^{-\frac{y^2}{2}} g_{m_1, \cdots, m_k}(y)
    (b+1 - y)^{m_{k+1}} \, \diff y
}{
    \displaystyle \int_{-\infty}^b \mathrm{e}^{-\frac{y^2}{2}} g_{m_1, \cdots, m_k}(y)
    (b - y)^{m_{k+1}} \, \diff y
} \\
\le\ & 2^{m_{k+1} - 1} \frac{
    \displaystyle \int_{-\infty}^b \mathrm{e}^{-\frac{y^2}{2}} g_{m_1, \cdots, m_k}(y)
    \bigl(1 + (b - y)^{m_{k+1}}\bigr) \, \diff y
}{
    \displaystyle \int_{-\infty}^b \mathrm{e}^{-\frac{y^2}{2}} g_{m_1, \cdots, m_k}(y)
    (b - y)^{m_{k+1}} \, \diff y
} \\
& + \frac{
    \displaystyle \int_0^1 \mathrm{e}^{-\frac{(b+1 - y)^2}{2}} g_{m_1, \cdots, m_k}(b+1 - y)
    y^{m_{k+1}} \, \diff y
}{
    \displaystyle \int_0^\infty \mathrm{e}^{-\frac{(b - y)^2}{2}} g_{m_1, \cdots, m_k}(b - y)
    y^{m_{k+1}} \, \diff y
} \\
\le\ & 2^{m_{k+1} - 1} + 2^{m_{k+1} - 1} 
\frac{
    \displaystyle \int_{-\infty}^b \mathrm{e}^{-\frac{y^2}{2}} g_{m_1, \cdots, m_k}(y) \, \diff y
}{
    \displaystyle \int_{-\infty}^{-(b)_- - 1} \mathrm{e}^{-\frac{y^2}{2}} g_{m_1, \cdots, m_k}(y) \, \diff y
} \\
& + C_k \frac{
    \displaystyle \int_0^1 \mathrm{e}^{-\frac{(b+1 - y)^2}{2}} (1 + \mathrm{e}^{-k(b - y)}) g_{m_1, \cdots, m_k}(b - y)
    y^{m_{k+1}} \, \diff y
}{
    \displaystyle \int_0^\infty \mathrm{e}^{-\frac{(b - y)^2}{2}} g_{m_1, \cdots, m_{k-1}}(b - y)
    y^{m_{k+1}} \, \diff y
}\\
&\displaystyle\le 2^{m_{k+1}-1} + 2^{m_{k+1}-1}\frac{\left(\int^b_{-(b)_- -1}+\int^{-(b)_- -1}_{-\infty}\right)\mathrm{e}^{-\frac{(y)^2}{2}} g_{m_1, \cdots, m_{k}}(y)\diff y}{\int^{-(b)_- -1}_{-\infty}\mathrm{e}^{-\frac{(y)^2}{2}} g_{m_1, \cdots, m_{k}}(y)\diff y}\\
         &\displaystyle +C_k\cdot\frac{\mathrm{e}^{k+1}(1+\mathrm{e}^{-kb}) \mathrm{e}^{-b}\int^1_{0}\mathrm{e}^{by-\frac{y^2}{2}}g_{m_1, \cdots, m_{k}}(b-y)(y)^{m_{k+1}}\diff y}{\int^\infty_{0}\mathrm{e}^{by-\frac{y^2}{2}} g_{m_1, \cdots, m_{k-1}}(b-y)(y)^{m_{k+1}}\diff y}\\
         &\displaystyle\le 2^{m_{k+1}} +C_k\mathrm{e}^{k+1-b} (1+\mathrm{e}^{-kb})\\
         &\displaystyle+\mathbf{1}_{b\le 0}2^{m_{k+1}-1}C_k\mathrm{e}^{k-b} (1+\mathrm{e}^{-k(b-1)})+\mathbf{1}_{b\ge0}\frac{2^{m_{k+1}-1}\int^b_{ -1}\mathrm{e}^{-\frac{(y)^2}{2}} g_{m_1, \cdots, m_{k}}(y)\diff y}{\int^{-1}_{-\infty}\mathrm{e}^{-\frac{(y)^2}{2}} g_{m_1, \cdots, m_{k}}(y)\diff y}\\
         &\displaystyle\le 2^{m_{k+1}} + C_k\mathrm{e}^{k+1-b} (1+\mathrm{e}^{-kb})+ \mathbf{1}_{b\le 0}2^{m_{k+1}-1}C_k\mathrm{e}^{k-b} (1+\mathrm{e}^{-k(b-1)})\\
         &+ 2^{m_{k+1}-1}\mathbf{1}_{b\ge0}\frac{\int^\infty_{ -\infty}\mathrm{e}^{-\frac{(y)^2}{2}} g_{m_1, \cdots, m_{k}}(y)\diff y}{\int^{-1}_{-\infty}\mathrm{e}^{-\frac{(y)^2}{2}} g_{m_1, \cdots, m_{k}}(y)\diff y}\,. 
\end{align*}

    Now observe that we can estimate
    \begin{equation*}
    \begin{array}{cc}
        &\displaystyle\int_{\R}\mathrm{e}^{-\frac{(y)^2}{2}} g_{m_1, \cdots, m_{k}}(y)\diff y\\
        & = \displaystyle \int_{\R}\int_{x_1\le x_2\le \cdots\le x_k\le y}\mathrm{e}^{-\frac{y^2}{2}}\prod_{i=1}^{k}\mathrm{e}^{-\frac{x_i^2}{2}}\cdot(x_2-x_1)^{m_1}\cdot \cdots \cdot (y-x_k)^{m_k}\diff x_1\cdots \diff x_{k}\diff y\\
        & \le \displaystyle \int_{\R^{k+1}}\mathrm{e}^{-\frac{y^2}{2}}\prod_{i=1}^{k}\mathrm{e}^{-\frac{x_i^2}{2}}\cdot|x_2-x_1|^{m_1}\cdot \cdots \cdot |y-x_k|^{m_k}\diff x_1\cdots \diff x_{k}\diff y\\
        & \le 2^{\sum_{i=1}^k m_i }\displaystyle \int_{\R^{k+1}}\prod_{i=1}^{k+1}\mathrm{e}^{-\frac{x_i^2}{2}}\cdot\big(|x_2|^{m_1}+|x_1|^{m_1}\big)\cdot \cdots \cdot \big(|x_{k+1}|^{m_k}+|x_k|^{m_k}\big)\diff x_1\cdots \diff x_{k+1}\\
        & \le 2^{\sum_{i=1}^k m_i + k}\mathrm{e}^{d\sum_{i=1}^k m_i\log m_i}\\
    \end{array}
    \end{equation*}
    for a universal constant $d>0$ using independence and the asymptotics of moments of a gaussian random variable. Furthermore, repeated applications of Lemma \ref{lemma: integral estimate} give the lower bound 
    \begin{equation*}
    \begin{array}{cc}
          &C_{k, \underline{m}}\le \displaystyle\int^{-1}_{-\infty}g_{m_1, \cdots, m_{k}}(y) \diff y\,,
    \end{array}
    \end{equation*}
for $C_{k, \underline{m}} = c^k\prod_{i=1}^k\frac{1}{i^{m_{i}}}$ for some universal constant $c>0$. Combining the above, we arrive at
\begin{equation*}
\begin{array}{cc}
     & \frac{g_{m_1, \cdots, m_{k+1}}(b+1)}{g_{m_1, \cdots, m_{k+1}}(b)}\le 2^{m_{k+1}} + \mathbf{1}_{b\le 0}2^{m_{k+1}-1}C_k\mathrm{e}^{k-b} (1+\mathrm{e}^{-k(b-1)})\\
     & +C_k\mathrm{e}^{k+1-b} (1+\mathrm{e}^{-kb}) + 2^{m_{k+1}-1}\mathbf{1}_{b\ge 0}c^k \cdot \mathrm{e}^{d\sum_{i=1}^k m_i\log m_i} \\
     & \le C_{k+1}(1+\mathrm{e}^{-(k+1)b})
\end{array}
\end{equation*}
    for universal constants $c,d>0$ where $C_{k+1} = 2^{m_{k+1}}(C_{k}\mathrm{e}^{k+1}\lor \mathrm{e}^{d\sum_{i=1}^k m_i\log m_i})$. This concludes the induction.
\end{proof}

\end{document}
